\paragraph*{In short}

\begin{table}[htbp]
\centering\small
\begin{tabularx}{\linewidth}{@{}lXX@{}}
\toprule
\textbf{Pattern} & \textbf{Primitives used} & \textbf{When to use} \\
\midrule
\texttt{swap} double buffering & \texttt{dma.any}, \texttt{swap} & Single thread group interleaving load and compute \\
Event pipeline & \co@inthreads.async@, events & Separate producer and consumer groups \\
\bottomrule
\end{tabularx}
\end{table}

Handle rotation is enough when one group controls both the transfer schedule and the computation. Separate producer and consumer regions need events because neither can infer the other's progress from its own instruction order. More buffering can give them room to overlap, but it cannot repair a missing dependency.

The same synchronization patterns can be combined with the richer data-movement primitives described in this guide.
